\item Un terremoto bajo el mar o un deslizamiento de Tierra puede producir una onda oceánica de corta duración transportando gran energía, llamada tsunami. Cuando su longitud de onda es grande comparada con la profundidad del océano $d$, la rapidez de una onda acuática está dada aproximadamente por $v = \sqrt{gd}$. Suponga que un terremoto ocurre a lo largo de la frontera de una placa tectónica que corre de norte a sur y produce una cresta de onda tsunami recta que en todas partes se mueve hacia el oeste.
\begin{enumerate}
    \item ¿Qué cantidad física puede considerarse constante en el movimiento de cualquier cresta de onda?
    \item Explique por qué la amplitud de la onda se incrementa conforme la onda se aproxima a la orilla.
    \item Si la onda tiene una amplitud de 1.80 m cuando su rapidez es de 200 m/s, ¿cuál será su amplitud donde el agua tenga 9.00 m de profundidad?
    \item Explique por qué se debe esperar que la amplitud sea mayor en la orilla, pero no se puede predecir significativamente mediante su modelo.
\end{enumerate}